\section{Specific Replies to Each Reviewer}
\label{sec:specific-replies}

\subsection{Reply to Reviewer-B}

\noindent\textbf{Suggestion B1:} \textit{In several places you write ``hide latency''; maybe it will be better to write ``reduce latency''.} 
\\

\noindent Thank you for the wording suggestion. We have updated all instances of ``hide latency'' and related wording to ``reduce amortized latency'' in the revised manuscript.
\\

\noindent\textbf{Suggestion B2:} \textit{In Ref.~\cite{gonzalez2012powergraph}, \{PowerGraph\} should be PowerGraph and \{Graph-Parallel\} should be Graph-Parallel.}
\\

\noindent Thank you for pointing out this typo. We have corrected Ref.~\cite{gonzalez2012powergraph} by removing the rendered braces around PowerGraph and Graph-Parallel.
\\

\subsection{Reply to Reviewer-C}

\noindent\textbf{Comment C1:} \textit{Unbalanced hardware tiers confound the CPU-GPU suitability claim. The CPU baselines are evaluated on an Intel Core i9-14900HX system, while the GPU baselines run on a separate machine with an NVIDIA Tesla A100. This comparison mixes a mobile-class CPU platform with a data-center-class GPU platform, making it difficult to attribute observed differences to architectural suitability rather than system-class effects (e.g., sustained power limits and, especially, host memory subsystem/bandwidth).}
\\

\noindent Thank you for raising this concern about comparison fairness. In the revised experiments, we moved both the CPU and GPU evaluations to server-class AWS platforms to avoid comparing a mobile-class CPU system against a data-center GPU system. The updated platform choice is summarized in Subsection~\ref{subsec:rev-fair-platform} of Section~\ref{sec:major-revisions}.
\\

\noindent\textbf{Comment C2:} \textit{In addition, the CPU configuration description appears ambiguous (on page 11) (the i9-14900HX is a mobile HX-series part and is commonly specified as 24 cores/32 threads), so clarifying cores vs.\ threads and the threading/pinning policy is important for reproducibility.}
\\

\noindent Thank you for pointing out this ambiguity. We revised the experimental setup to explicitly distinguish CPU cores from hardware threads and to include the relevant CPU architectural details for reproducibility. These clarifications are described in Subsection~\ref{subsec:rev-cores-threads} and Subsection~\ref{subsec:rev-exp-setup} of Section~\ref{sec:major-revisions}.
\\

\subsection{Reply to Reviewer-D}

\noindent\textbf{Comment D1:} \textit{It would be good to explain why the ratio is small (close to 1) for some of the larger graphs in Table 5 (column 6) if enough parallel work is available under XB++?}
\\

\noindent This is an insightful point. After moving the experiments from the A100 platform to the RTX PRO 6000 Blackwell platform, this issue no longer appears: XB++ now achieves a higher instruction issue rate than BFS-Gunrock on all evaluated graphs. The revised results suggest that the earlier near-1 ratios were partly due to limited available SM-level execution resources on the original GPU platform rather than a lack of parallel work in XB++. We address this ratio issue in Subsection~\ref{subsec:rev-inst-ratio} of Section~\ref{sec:major-revisions}.
\\

\noindent\textbf{Comment D2:} \textit{It would be nice to comment on the variability in the improvement of the speedup across different graphs in Table 2.}
\\

\noindent We appreciate the reviewer highlighting this variability. To explain the differences across graphs, we added a new metric, Reuse-Ratio, which measures how effectively the alternating-tree reuse mechanism preserves useful tree nodes instead of resetting them. This metric helps connect the observed speedups in Table~2 to the internal behavior of the reuse optimization. The added explanation of speedup variability is given in Subsection~\ref{subsec:rev-reuse-variability} of Section~\ref{sec:major-revisions}.
\\

\noindent\textbf{Comment D3:} \textit{It would be nice to expand a bit on how SIMT execution helps to keep up massive concurrency (I am not very familiar with this area).}
\\

\noindent This is necessary context for interpreting the GPU performance results. In the revision, we expanded the discussion to explain how resident warps help the GPU maintain high concurrency when XB++ exposes enough fine-grained edge-level work. The added discussion is summarized in Subsection~\ref{subsec:rev-simt-concurrency} of Section~\ref{sec:major-revisions}.
\\

\noindent\textbf{Comment D4:} \textit{Can you please explain small working set a bit more on Page 17: is it related to the amount of memory needed to store required data in cache?}
\\

\noindent This is an important clarification. We added an explicit definition of working set in the data-movement discussion of the profiling methodology. The corresponding revision is summarized in Subsection~\ref{subsec:rev-working-set} of Section~\ref{sec:major-revisions}.
\\

\noindent\textbf{Comment D5:} \textit{Section 2.4 is not very clear: How are the 3 procedures performed parallely? It would be nice to expand on the claim in the introduction that the usual parallel algorithm throws away alternating paths instead of reusing it.}
\\

\noindent We agree that the parallel execution of the three procedures needed a clearer explanation. We added a paragraph in Section~2.4 describing how the current checking set is split across worker threads and how each phase processes its assigned frontier vertices. The added paragraph is marked in \hyperref[para:xblossom-three-parallel-procedures]{this paragraph}.
\\

\noindent\textbf{Comment D6:} \textit{The acronyms CPU and GPU must be expanded the first time it is used in the paper -- same goes for SIMT and LLC.}
\\

\noindent Thank you for catching these acronym issues. We expanded CPU and GPU at their first prose occurrence in \hyperref[para:first-cpu-gpu]{the abstract} and also clarified them in the experimental setup at \hyperref[para:exp-setup-cpu]{the CPU platform description} and \hyperref[para:exp-setup-gpu]{the GPU platform description}. We also expanded LLC at \hyperref[para:first-llc]{its first body-text occurrence} and SIMT at \hyperref[para:first-simt]{its first evaluation occurrence}.
\\

\noindent\textbf{Comment D7:} \textit{Figure~\ref{plot:load-balance} discussion: It would be nice to remind the reader that having blossoms reduces the effectiveness of node-level parallelism as discussed in Subsection~\ref{subsubsec:exp_load_balancing}.}
\\

\noindent This connection helps interpret the load-balancing results. We added the explanation that blossom-heavy workloads involve sequential path-table operations, including tracing endpoint-to-root paths to identify the blossom base and materializing the odd cycle from path segments. The added discussion is marked in \hyperref[para:blossom-node-level-limitation]{this paragraph}.
\\

\noindent\textbf{Comment D8:} \textit{Section~\ref{sec:intensive-experiments} discusses BFS-Ligra and BFS-Gunrock but they are compared only in Section~\ref{sec:arch_suitability}. This discussion could possibly be moved to the start of Section~\ref{sec:arch_suitability}.}
\\

\noindent Thank you for the organization suggestion. The moved discussion is now in Subsection~\ref{subsec:baselines}.
\\

\subsection{Reply to Reviewer-E}

\noindent\textbf{Contribution Concern E1:} \textit{The proposed profiling methodology is not new. It is using well-known common metrics and popular tools. It is a good case study but I think the profiling tehcniques and metrics you proposed are standard procedures. Additionally, the selection of the LLC and memory efficiency needs more scientific justification. How about other metrics? .e.g, LLC occupancy, effective DRAM bandwidth? I think a more careful consideration of the metric space exploration and narrowing down to a limited set of metrics would strength the technical depth of this part.
You claim your contribution in page 3 ``Our performance analysis methodology can be applied to general parallel-processing applications'': I am assuming you are referring to the profiling part. The profiling technique is standard, so I agree that this can be applied to general applications. But this cannot be listed as a contribution. I would recommend using this bullet point to talk about the takeaway messages you obtained from your profiling case study.}
\\

\noindent Thanks for your constructive suggestions.
To justify our selected metrics, we added a paragraph before introducing our measures to explain why commonly reported peak metrics, including peak FP64/FP32/FP16/integer throughput and peak memory bandwidth, do not capture the execution pattern of traversal-based graph workloads.
This metric-selection rationale has been added in \hyperref[para:metric-selection-justification]{this paragraph}.

\noindent We also revised the contribution statement to focus on the architectural-suitability takeaway distilled from the case study. The revised contribution bullet is marked in \hyperref[item:takeaway-architecture-suitability]{this contribution item}.
\\

\noindent\textbf{Contribution Concern E2:} \textit{All the evaluated graphs are sparse. Given that many of the real-world graph networks are sparse, your contribution is still valuable. However, please either tone down your contribution from general graphs to sparse graphs, or you should include dense graph benchmarks in your evaluation.}
\\

\noindent Thank you for this helpful scope clarification.
We revised the dataset description to state explicitly that this paper focuses on real-world sparse graph datasets.
The added clarification is marked in \hyperref[para:real-world-sparse-datasets]{this paragraph}.
\\

\noindent\textbf{Contribution Concern E3:} \textit{The whole paper covers two types of graph algorithms(applications): BFS and blossom, but they are loosely connected. The argument from the authors is that these two algorithms are both graph traversal with frontier-expansion paradigm. Then how about other traversal algorithms? e.g., SSSP?}
\\

\noindent We appreciate this suggestion to strengthen the generality of our performance-analysis approach beyond BFS and blossom. We have already applied it to two additional traversal-based graph algorithms, SSSP and MSSP. Their instruction execution rates are reported in Table~\ref{tab:pe}, and their CPU cache behavior and GPU memory-bandwidth results are reported in Table~\ref{tab:memory}, both in Section~\ref{sec:arch_suitability}.
\\

\noindent\textbf{Evaluation Detail E1:} \textit{What framework/lib are you using for evaluation? For CPU, are you using openMP, or pthread? For the two baselines Ligra and Gunrock, what are they using? These things should be included in the paper to judge that you are doing a fair comparison.}
\\

\noindent Thank you for the question. The added framework and environment details are summarized in Subsection~\ref{subsec:rev-exp-setup} of Section~\ref{sec:major-revisions}. In the main paper, XB-Pro and XB++ parallelization details are given in Subsubsection~\ref{subsubsec:eval-exp-env}, while the Ligra and Gunrock baseline details are given in Subsection~\ref{subsec:baselines}.
\\

\noindent\textbf{Evaluation Detail E2:} \textit{The scalability test is missing. When varying the number of threads, how does the performance trend look like?}
\\

\noindent Thank you for pointing out this. We added a scalability test that varies the number of CPU threads for XB-Pro and the number of enabled GPU SMs for XB++.
The revision is summarized in Subsection~\ref{subsec:rev-scalability} of Section~\ref{sec:major-revisions}, with full details in Subsection~\ref{subsec:scalability_test}.
\\

\noindent\textbf{Evaluation Detail E3:} \textit{The CPU platform uses i9-14900HX, which has 24 cores, not 32 cores. It is 8 performance cores x 2 (SMT) + 16 Efficient cores = 32 threads when SMT (hyperthreading) is enabled.}
\\

\noindent Thank you for the correction. We revised the core/thread description, as documented in Subsection~\ref{subsec:rev-cores-threads} of Section~\ref{sec:major-revisions}.
\\

\noindent\textbf{Evaluation Detail E4:} \textit{Confidence intervals missing.}
\\

\noindent Thank you for the suggestion. We added confidence intervals for all experiments to strengthen the reliability of the reported results, as described in Subsection~\ref{subsec:rev-confidence-intervals} of Section~\ref{sec:major-revisions}.
\\

\noindent\textbf{Evaluation Detail E5:} \textit{Lack of details of execution phases: how much portion is sequential, how much can be run in parallel? This would help in understanding the parallel efficiency of your method. Currently, for some workloads, e.g., Wikipedia and Youtube, the CPU$\rightarrow$GPU performance speedup can go beyond 10x, but for some workloads like GPlus, HiggsNets, the CPU$\rightarrow$GPU performance speedup is around 4x (table 2). This does not match the growth of thread count from CPU to GPU. Clarifying the parallel vs sequential portion might help with this concern.}
\\

\noindent Thank you for raising this concern. We added a description of execution phases, and the revised execution-phase discussion is summarized in Subsection~\ref{subsec:rev-execution-phases} of Section~\ref{sec:major-revisions}.
\\

\noindent\textbf{Presentation Suggestion E1:} \textit{I would suggest the authors to turn some tables into plots (e.g., bar charts). The current tables are hard to compare.}
\\

\noindent Thank you for the suggestion. We converted the load-balancing runtime table into a bar-plot figure to make the node-level and edge-level comparisons easier to read across datasets. The revised plot is shown in Figure~\ref{plot:load-balance}.
\\

\noindent\textbf{Presentation Suggestion E2:} \textit{From the title, it looks like the paper is a characterization paper. However, sec 3 presents your optimization, sec 4 presents evaluation, and sec 5 presents a profiling-based case study. Then the paper reads like a mix of proposing new performance optimizations and characterization. I would suggest the authors to think about whether this is a characterization paper, or a perf optimization paper.
Currently as a characterization paper, the paper does not have enough interesting and thoughtful insights. Section 5.2.3 provides some takeaway messges: 1) XB++ runs better on GPU - This is new. 2) BFS is better on CPU - This is well known, as BFS is a challenging algorithm on GPUs due to the irregular graph data structure. The takeaway messages can be strengthened by generalizing and discussing more graph algorithms, or discuss future blossom/BFS hw-algorithm co-design.
As a perf optimization paper, the algorithm, implementation, and profiling contributions are kind of thin. You are applying some well-known techniques to a new problem. But most of them are engineering efforts. Even the algorithm (reuse), it is inspired/motivated by bipartite graphs [3].}
\\

\noindent Thank you for this suggestion. To strengthen the takeaway messages and broaden the characterization beyond BFS and blossom, we added two more traversal-based graph algorithms, SSSP and MSSP, to the profiling study. The expanded results are reported in Table~\ref{tab:pe} for instruction execution rates and Table~\ref{tab:memory} for CPU cache behavior and GPU memory-bandwidth behavior in Section~\ref{sec:arch_suitability}.
\\
